\section{Sharp obstructions and the existing-theorem boundary}
\label{sec:tpc33-obstructions-literature}

The preceding theorems create an average route around pointwise affine
cancellation.  They do not make a prescribed direction small.  We record
three exact obstructions before comparing the physical requirement with
existing Chowla--Elliott technology.

\subsection{Three sharp interface obstructions}
\label{sec:tpc33-sharp-obstructions}

\begin{proposition}[A fixed direction contains two-point Chowla]
\label{prop:tpc33-chowla-subcase}
The family \(sU-aD=h\) contains the shift-one M\"obius correlation.  In
particular, take
\begin{equation}
 a=s=h=1,\qquad x=1,\qquad y=0.
 \label{eq:tpc33-shift-one-parameters}
\end{equation}
Then
\begin{equation}
 D_t=t,\qquad U_t=t+1,\qquad
 \mu(D_t)\mu(U_t)=\mu(t)\mu(t+1).
 \label{eq:tpc33-shift-one-correlation}
\end{equation}
Thus an unweighted \(o(T)\) theorem uniform over every individual affine
direction would imply the still open two-point M\"obius Chowla statement.
\end{proposition}

\begin{proof}
The B\'ezout relation is \(sx-ay=1\).  Substitution in
\eqref{eq:tpc33-canonical-affine-forms} gives
\eqref{eq:tpc33-shift-one-correlation}.
\end{proof}

The next example shows why separate Fourier pseudorandomness of the two
legs cannot control a prescribed pairing.

\begin{proposition}[Fourier-flat legs with a coherent line]
\label{prop:tpc33-coherent-line-counterexample}
Let \(q\) be an odd prime, fix \(a_0,s_0,h\in\F_q^\times\), and let
\(\chi\) be the quadratic character with \(\chi(0)=0\).  Define
\begin{equation}
 f(D)=\chi(D),
 \qquad
 g(U)=\chi\!\left(a_0^{-1}(s_0U-h)\right).
 \label{eq:tpc33-legendre-coherent-functions}
\end{equation}
Then \(\bar f=\bar g=0\), and every nonzero additive Fourier coefficient
of either function has modulus \(\sqrt q\).  Nevertheless,
\begin{equation}
 \boxed{
 \cC_{a_0,s_0}^{(h)}(f,g)=q-1.}
 \label{eq:tpc33-legendre-coherent-line}
\end{equation}
\end{proposition}

\begin{proof}
Both functions are affine transforms of the quadratic character, so their
means vanish and their nonzero Fourier coefficients are quadratic Gauss
sums of modulus \(\sqrt q\).  On the selected line
\(U=s_0^{-1}(a_0D+h)\),
\[
 g(U)=\chi\!\left(a_0^{-1}(a_0D+h-h)\right)=\chi(D)=f(D).
\]
Hence the correlation is \(\sum_D\chi(D)^2=q-1\).
\end{proof}

This counterexample concerns an abstract two-leg Fourier interface.  It
does not assert that the M\"obius function possesses the same coherent
structure; it proves that one-leg Fourier bounds alone cannot rule it out.

\begin{proposition}[The single-modulus no-wrap conflict]
\label{prop:tpc33-single-modulus-conflict}
On the high-\(\beta\) physical packet, a single prime modulus cannot both
turn the modular determinant condition into an exact integer equation on
the ambient affine box and provide complete orbit-slope coverage.
\end{proposition}

\begin{proof}
On a projected physical layer, the slopes are \(s\) and
\(a=\ell jv\), while the \(\widetilde d\)-interval has ambient length
\(\asymp D/v\).  Therefore the determinant expression varies on the scale
\begin{equation}
 a\frac Dv=\ell jD\asymp QJ\asymp X.
 \label{eq:tpc33-ambient-determinant-span}
\end{equation}
An elementary no-wrap lift from
\(sU-aD\equiv h\pmod q\) to equality throughout this ambient box requires
\(q\gg X\).  In contrast, complete coverage by the orbit variable
\(j\) requires \(q\lesssim J\); otherwise the rms factor in
\eqref{eq:tpc33-incomplete-rms-loss} contains \((q/J)^{1/2}\).  At the
physical endpoint \(J=X^{133/400+o(1)}\ll X\), these requirements are
incompatible.
\end{proof}

The proposition rules out only the most direct one-modulus transplantation
of \cref{thm:tpc33-projective-plancherel}.  It does not rule out a
multi-modulus circle method, a delta symbol, a spectral decomposition, or a
direct integer large-sieve inequality for the restricted slope set.

\subsection{What the proved affine-correlation theorems supply}

Several strong results control related averages.  Their hypotheses do not
coincide with the physical column.  The comparison in
\cref{tab:tpc33-literature-boundary} uses only the theorem forms stated in
the cited primary sources.  The earlier TPC-23 paper already records the
fixed-form logarithmic bridge in the present reduction
\citep{WangTPC23}; it is not claimed again as a new theorem here.

\begin{table}[t]
\centering
\caption{Existing affine-correlation inputs and the missing physical
interfaces.}
\label{tab:tpc33-literature-boundary}
\small
\begin{tabular}{@{}>{\raggedright\arraybackslash}p{0.20\textwidth}
>{\raggedright\arraybackslash}p{0.20\textwidth}
>{\raggedright\arraybackslash}p{0.23\textwidth}
>{\raggedright\arraybackslash}p{0.28\textwidth}@{}}
\toprule
Source & Coefficients & Averaging & Boundary for TPC-33 \\
\midrule
Tao, Cor.~1.5 \citep{Tao2016LogChowla}
& fixed nonproportional affine forms
& logarithmic window \(x/\omega(x)<n\le x\), \(\omega(x)\to\infty\)
& applies to each frozen M\"obius pair; no stated uniformity for
polynomially growing coefficients and not a fixed dyadic window \\
Matom\"aki--Radziwi\l\l--Tao, Thm.~1.6
\citep{MatomakiRadziwillTao2015}
& coefficients bounded in terms of \(A\), with an \(A^2\) loss
& unweighted average plus a free additive-shift average
& applies also to M\"obius, but the displayed bound yields no saving when
\(A=X^\theta\); the
free shift is not the determinant-constrained family \\
Frantzikinakis--Host, Thm.~2.6 \citep{FrantzikinakisHost2017}
& fixed independent homogeneous linear forms
& a full two-dimensional box
& gives mean zero for aperiodic M\"obius, but not on the one-dimensional
affine slice or for growing forms \\
Matom\"aki et al., Cor.~1.11
\citep{MatomakiEtAl2023HigherUniformity}
& fixed distinct dilation coefficients
& average over a common dilation
& strong precedent for extra slope averaging; the displayed corollary is
for Liouville and does not select one prescribed dilation \\
\bottomrule
\end{tabular}
\end{table}

Two further nearby boundaries are worth recording.  Ter\"av\"ainen proves
a fixed gap from the trivial bound for certain roots-of-unity-valued
multiplicative functions on fixed nonproportional affine forms
\citep[Thm.~2.6]{Teravainen2024PolynomialArguments}; this is not a mean-zero theorem,
and the roots-of-unity hypothesis does not directly permit \(\mu\), which
also takes the value zero.  Pilatte proves a quantitative logarithmic
shift-one result for the Liouville function
\citep[Thm.~1.1]{Pilatte2025LogChowla}; it is a preprint result for a particular fixed
pair, not a general growing-coefficient dyadic M\"obius theorem.

In summary, no cited theorem simultaneously supplies:
\begin{enumerate}[label=\textup{(\roman*)}]
 \item polynomially growing affine coefficients;
 \item cancellation on every fixed dyadic shell;
 \item averaging along the fixed-determinant constrained family;
 \item the actual joint row mask and physical weights; and
 \item a power ledger strong enough for \eqref{eq:tpc33-Q3-energy-input}.
\end{enumerate}
These requirements remain an additional analytic input.  The projective
identity in this paper identifies one viable averaged geometry, but does
not claim that existing Chowla--Elliott theorems already close it.
The distinction between sign-sensitive prime detection and sign-blind
sieve information is the classical parity boundary
\citep{HeathBrown1982Parity}.
